\section{Semantic Session Typing for Deadlock Freedom}
\label{sec:logicalRelation}

The semantic typing judgment $\senv \proves \proc : \ty$ of this section
fulfils three goals.
\begin{enumerate}[(1)]
  \item It ensures deadlock freedom for closed processes (``adequacy'');
  \item
    It validates all rules of the static system of
    \cref{sec:syntacticTyping}, which gives the ``fundamental property''
    that
    $
      \senv \types \proc : \ty
      \implies
      \senv \proves \proc : \ty
    $;
  \item It contains safe but untypable processes.
\end{enumerate}


\subsection{Worlds and Readiness}
\label{sec:worlds}

The relation is indexed by a \emph{world}, which carries the order the
points of a running program stand in.

\begin{definition}[World \rocq{world}, \rocq{world_wf}, \rocq{strict}]
\label{def:world}
  A world $\wrld$ is a finite set of names, its \emph{support}, together
  with two relations $\fEq{}{}$ and $\fLt{}{}$ on points. It is
  \emph{well formed} when $\fEq{}{}$ is an equivalence closed under
  successor and slot congruence, $\fLt{}{}$ is transitive, closed under
  $\fEq{}{}$ on both sides and contains program order
  $\fLt{\pntA}{\poff{\pntA}{1}}$, every support name is a point
  constant, and outside the support only reflexivity and program order
  hold. It is \emph{strict} when $\fLt{}{}$ is irreflexive.
  $\wrld$ \emph{knows} $\pntA$ when $\rootof(\pntA)$ is in its support
  (\rocq{pt_known}).
\end{definition}

\begin{definition}[Extension \rocq{wext}, \rocq{wext_fresh}]
\label{def:wext}
  $\wrld \wextsym \wrld'$ holds when $\wrld'$ is strict and contains the
  support and both relations of $\wrld$.
  $\wrld \wextfresh \wrld'$ \emph{placing $N$} holds in addition when
  $N$ is disjoint from the support of $\wrld$, the support of $\wrld'$
  is that of $\wrld$ with $N$, and $\wrld'$ adds no fact between points
  whose roots avoid $N$.
\end{definition}

Only $\tNew$ extends the world, and it extends it freshly. The two
results below are what makes a composed run strict.

\begin{theorem}[Amalgamation \rocq{amalgamation}]
\label{thm:amalgamation}
  For strict $\wrld_1$ and $\wrld_2$, the join
  $\wrld_1 \sqcup \wrld_2$ is strict if and only if the join of their
  restrictions to the names both support is strict.
\end{theorem}

\begin{theorem}[Fresh extension \rocq{fresh_extension}]
\label{thm:fresh-extension}
  If $\wrld_1 \wextfresh \wrld_1'$ placing $N$, the support of
  $\wrld_2$ avoids $N$, $\wrld_2$ is strict and
  $\wrld_1 \sqcup \wrld_2$ is strict, then
  $\wrld_1' \sqcup \wrld_2$ is strict.
\end{theorem}

The step clauses of the relation are gated by \emph{readiness}. A term
or process with no internal step left must either be done or offer some
of the entries it holds, and the offered entries must lie strictly below
everything it holds but does not offer. Entries are compared, not names:
a thread holding both ends of $\ch$ and blocking on one of them holds
the other entry unoffered and would need
$\fLt{\heads(\ch)}{\heads(\ch)}$.

\begin{definition}[Readiness \rocq{ready_t}, \rocq{ready_p}]
\label{def:ready}
  $\readyop_{\wrld}(\tenv \mid \frm, \tm)$ holds when, provided $\tm$
  has no $\beta$, $\lNew{\ch}$ or $\lSpawnL{\val}$ step,
  \begin{enumerate*}[(i)]
    \item $\tm$ is a value if it offers nothing; and
    \item $\tm$ has a forwarding redex, or every entry it offers lies in
      $\wrld$ strictly below every entry of $\tenv$ it does not offer
      and strictly below every point of $\frm$.
  \end{enumerate*}
  $\readyop_{\wrld}(\senv, \proc)$ holds when, provided $\proc$ has no
  $\tau$ step, $\proc$ has a forwarding redex, or $\proc$ is terminated
  if it offers nothing and every entry it holds and does not offer lies
  strictly above some entry it offers.
\end{definition}

A forwarding never blocks, so a term or process at a forwarding redex
compares nothing.

\subsection{A Logical Relation for \OurLang}
\label{sec:logicalRelation:relation}

Three relations express the typing constraints
(\cref{fig:logrel-rules-value,fig:logrel-rules-term,fig:logrel-rules-proc}):
$\lrVW{\wrld}{n}{\tenv \types \ty}$ for values,
$\lrEW{\wrld}{n}{\tenv \types \ty}{\frm}$ for terms, and
$\lrPW{\wrld}{n}{\senv \types \ty}$ for processes.
The index $n \in \Nat$ is a step index~\cite{Ahmed06} bounding the
number of transitions the relation examines.
A session environment $\senv$ maps each channel to one
$\sAt{\sty}{\pntA}$, and a term environment $\tenv$ maps a channel to
one or both of its ends, because a term may hold both
(\rocq{senv}, \rocq{tenv}). $\heads(\tenv)$ is the set of points of the
entries of $\tenv$, and $\tenv_1 \cdot \tenv_2$ splits $\tenv$ up to
$\fEq{}{}$ of the world (\rocq{tenv_comp}).

\begin{figure}[tp]
  \adjustfigure[]
  \rulesection*{Value logical relation $\val \in \lrVW{\wrld}{n}{\tenv \types \ty}$ (\rocq{Vdef})}
  \begin{proofrules}
    \infer*[lab=V-Unit]{}{
      \tUnit \in \lrVW{\wrld}{n}{\tEmpty \types \tyUnit}
    } \label{rule:V-Unit}
    \and
    \infer*[lab=V-Chan]{
      \sty_0 \approx \sty \\
      \wrld \types \fEq{\pntB}{\pntA}
    }{
      \ch \in \lrVW{\wrld}{n}{\ch : \sAt{\sty_0}{\pntB} \types \sAt{\sty}{\pntA}}
    } \label{rule:V-Chan}
    \and
    \infer*[lab=V-Const]{
      \cst : \ty_0 \text{ an instance of a signature under } \wrld \\
      \ty_0 \approx_{\wrld} \ty
    }{
      \cst \in \lrVW{\wrld}{n}{\tEmpty \types \ty}
    } \label{rule:V-Const}
    \and
    \infer*[lab=V-Pair]{
      \val_1 \in \lrVW{\wrld}{n}{\tenv_1 \types \ty} \\
      \val_2 \in \lrVW{\wrld}{n}{\tenv_2 \types \tyB}
    }{
      \tPair{\val_1}{\val_2} \in \lrVW{\wrld}{n}{\tenv_1 \cdot \tenv_2 \types \lPair{\ty}{\tyB}}
    } \label{rule:V-Pair}
    \and
    \infer*[lab=V-Inj]{
      \val \in \lrVW{\wrld}{n}{\tenv \types \ty_{\asTy{i}}}
    }{
      \tInj{\asTm{i}}{\val} \in \lrVW{\wrld}{n}{\tenv \types \lSum{\ty_1}{\ty_2}}
    } \label{rule:V-Inj}
    \and
    \infer*[lab=V-Abs]{
      \fchans(\tm) = \dom(\tenv) \\
      \wrld \types \fLe{\lbd}{\heads(\tenv)}
      \\\\
      \forall \wrld' \wextsym \wrld,\ j < n,\ \tenv'',\
        \val \in \lrVW{\wrld'}{j}{\tenv' \types \ty},\
        \frm \text{ with } \wrld' \types \fLt{\ubd}{\frm}
      \\
      \text{and } \tenv'' = \tenv \cdot \tenv' \st
        \tSub{\tm}{\val}{\vr} \in \lrEW{\wrld'}{j}{\tenv'' \types \tyB}{\frm}
    }{
      \tAbs{\vr}{\tm} \in \lrVW{\wrld}{n}{\tenv \types \ty \lolli{\ubd}{\lbd} \tyB}
    } \label{rule:V-Abs}
    \and
    \infer*[lab=V-Box]{
      \fchans(\tm) = \emptyset
      \\\\
      \forall \wrld' \wextsym \wrld,\ \pntvec \text{ known to } \wrld'
        \text{ with } \wrld' \types \tysub{\cstr}{\pntvec}{\pvivec} \st
      \\
      \tAbs{\tUnit}{\tSub{\tm}{\tBoxP{\vr}{\tm}}{\vr}} \in
        \lrVW{\wrld'}{n}{\tEmpty \types \tysub{(\tyUnit \lolli{\ubd}{\pTop} \ty)}{\pntvec}{\pvivec}}
    }{
      \tBoxP{\vr}{\tm} \in \lrVW{\wrld}{n}{\tEmpty \types \tyBox{\pvivec}{\cstr}{\tyUnit \lolli{\ubd}{\pTop} \ty}}
    } \label{rule:V-Box}
    \and
    \infer*[lab=V-Ex]{
      \wrld \types \fLe{\lbd}{\heads(\tenv)} \\
      \pntvec \text{ known to } \wrld \\
      \wrld \types \tysub{\cstr}{\pntvec}{\pvivec} \\
      \val \in \lrVW{\wrld}{n}{\tenv \types \tysub{\ty}{\pntvec}{\pvivec}}
    }{
      \val \in \lrVW{\wrld}{n}{\tenv \types \tyEx{\lbd}{\pvivec}{\cstr}{\ty}}
    } \label{rule:V-Ex}
  \end{proofrules}
  \rulesectionend
  \caption{The value logical relation}
  \Description{The seven rules of the value relation.}
  \label{fig:logrel-rules-value}
\end{figure}

\begin{figure}[tp]
  \adjustfigure[]
  \rulesection*{Term logical relation $\tm \in \lrEW{\wrld}{n}{\tenv \types \ty}{\frm}$ (\rocq{Edef})}
  \begin{proofrules}
    \infer*[lab=E-Zero]{}{
      \tm \in \lrEW{\wrld}{0}{\tenv \types \ty}{\frm}
    } \label{rule:E-Zero}
    \and
    \infer*[lab=E-Intro]{
      \irred \tm \implies \tm \in \lrVW{\wrld}{n+1}{\tenv \types \ty}
      \\
      \neg\irred \tm \implies \readyop_{\wrld}(\tenv \mid \frm, \tm)
      \\
      \forall \lambda, \tm'\st \tm \trans{\lambda} \tm' \implies
        \obl{\wrld ; \tenv \mid \frm}{n+1}{\tm \trans{\lambda} \tm' : \ty}
    }{
      \tm \in \lrEW{\wrld}{n+1}{\tenv \types \ty}{\frm}
    } \label{rule:E-Intro}
  \end{proofrules}
  \rulesection{Obligations of a term transition $\obl{\wrld ; \tenv \mid \frm}{n}{\tm \trans{\lambda} \tm' : \ty}$ (\rocq{Eobl})}
  \begin{proofrules}
    \infer*[lab=O-Beta]{
      \tm' \in \lrEW{\wrld}{n \cost{\lambda}}{\tenv \types \ty}{\frm}
    }{
      \obl{\wrld ; \tenv \mid \frm}{n}{\tm \trans{\beta} \tm' : \ty}
    } \label{rule:O-Beta}
    \and
    \infer*[lab=O-New]{
      \ch \notin \dom(\tenv) \implies
      \exists\, \pcst \text{ fresh for } \wrld, \tenv, \ty, \frm,\
        \sty,\ \wrld \wextfresh \wrld' \text{ placing } \pcst \st
      \\
      \tm' \in \lrEW{\wrld'}{n \cost{\lambda}}{\tenv , \ch : \sAt{\sty}{\pcst} \cdot \sAt{\dual{\sty}}{\pcst} \types \ty}{\frm}
    }{
      \obl{\wrld ; \tenv \mid \frm}{n}{\tm \trans{\lNew{\ch}} \tm' : \ty}
    } \label{rule:O-New}
    \and
    \infer*[lab=O-Spawn]{
      \val \in \lrVW{\wrld}{n}{\tenv_1 \types \tyUnit \lolli{\ubd}{\lbd} \tyUnit} \\
      \tm' \in \lrEW{\wrld}{n \cost{\lambda}}{\tenv_2 \types \ty}{\frm}
    }{
      \obl{\wrld ; \tenv_1 \cdot \tenv_2 \mid \frm}{n}{\tm \trans{\lSpawnL{\val}} \tm' : \ty}
    } \label{rule:O-Spawn}
    \and
    \infer*[lab=O-Recv]{
      \forall \wrld' \wextsym \wrld,\ j \le n \cost{\lambda},\ \tenv_2,\
        \val \in \lrVW{\wrld'}{j}{\tenv_2 \types \tysub{\tyB}{\pntA}{\pself}} \st
      \\
      \tm' \in \lrEW{\wrld'}{j}{(\tenv_1 , \ch : \sAt{\sty}{(\poff{\pntA}{1})}) \cdot \tenv_2 \types \ty}{\frm}
    }{
      \obl{\wrld ; \tenv_1 , \ch : \sAt{(\sRecv{\tyB}{\sty})}{\pntA} \mid \frm}{n}{\tm \trans{\lRecv{\ch}{\val}} \tm' : \ty}
    } \label{rule:O-Recv}
    \and
    \infer*[lab=O-Send]{
      \val \in \lrVW{\wrld}{n}{\tenv_2 \types \tysub{\tyB}{\pntA}{\pself}} \\
      \tm' \in \lrEW{\wrld}{n \cost{\lambda}}{\tenv_1 , \ch : \sAt{\sty}{(\poff{\pntA}{1})} \types \ty}{\frm}
    }{
      \obl{\wrld ; (\tenv_1 , \ch : \sAt{(\sSend{\tyB}{\sty})}{\pntA}) \cdot \tenv_2 \mid \frm}{n}{\tm \trans{\lSend{\ch}{\val}} \tm' : \ty}
    } \label{rule:O-Send}
    \and
    \infer*[lab=O-Close]{
      \tm' \in \lrEW{\wrld}{n \cost{\lambda}}{\tenv_1 \types \ty}{\frm}
    }{
      \obl{\wrld ; \tenv_1 , \ch : \sAt{\sTerm}{\pntA} \mid \frm}{n}{\tm \trans{\lTerm{\ch}} \tm' : \ty}
    } \label{rule:O-Close}
    \and
    \infer*[lab=O-Link]{
      \sty_{\chB} \isdual \sty_{\ch} \\
      \wrld \types \fEq{\pntA}{\pntB} \\
      \tm' \in \lrEW{\wrld}{n \cost{\lambda}}{\tenv_1 \types \ty}{\frm}
    }{
      \obl{\wrld ; \tenv_1 , \ch : \sAt{\sty_{\ch}}{\pntA} , \chB : \sAt{\sty_{\chB}}{\pntB} \mid \frm}{n}{}
      \\ \qquad {\tm \trans{\lLink{\ch}{\chB}} \tm' : \ty}
    } \label{rule:O-Link}
  \end{proofrules}
  \rulesectionend
  \caption{The term logical relation}
  \Description{The two rules of the term relation and the seven
    obligation rules of a term transition.}
  \label{fig:logrel-rules-term}
\end{figure}

A send, a close and a forwarding give a resource away, so they keep the
index. $n \cost{\lambda} = n$ for those three labels and
$n \cost{\lambda} = n - 1$ for every other (\rocq{tcost}). The
corresponding successor term is strictly smaller, which is what makes
\cref{fig:logrel-rules-term} a definition (\rocq{Efuel}).
\Cref{rule:V-Abs} reads the argument and the result at strictly smaller
indices only, so \ref{rule:V-Box} may use it at an instance of its own
type.

\begin{remark}
\label{rm:vabs}
  The side conditions on free channels of \ref{rule:V-Abs} and
  \ref{rule:V-Box} are needed because the LTS decides which channels a
  send extrudes by looking at the channels the sent value mentions
  syntactically (\cref{rm:sentc}). They are the only syntactic
  conditions in the relation, and they rule out a sent value that
  mentions channels in dead code which the continuation of the send
  still uses.
\end{remark}

\begin{figure}[tp]
  \adjustfigure[]
  \rulesection*{Process logical relation $\proc \in \lrPW{\wrld}{n}{\senv \types \ty}$ (\rocq{Pdef})}
  \begin{proofrules}
    \infer*[lab=P-Zero]{}{
      \proc \in \lrPW{\wrld}{0}{\senv \types \ty}
    } \label{rule:P-Zero}
    \and
    \infer*[lab=P-Intro]{
      \irred \proc \implies \proc = (\asPrOp{\emptyset}, \mset{\val,\tUnit,\dots,\tUnit})
        \text{ with } \val \in \lrVW{\wrld}{n+1}{\senv \types \ty}
      \\
      \neg\irred \proc \implies \readyop_{\wrld}(\senv, \proc)
      \\
      \forall \pi, \proc'\st \proc \trans{\pi} \proc' \implies
        \obl{\wrld ; \senv}{n+1}{\proc \trans{\pi} \proc' : \ty}
    }{
      \proc \in \lrPW{\wrld}{n+1}{\senv \types \ty}
    } \label{rule:P-Intro}
  \end{proofrules}
  \rulesection{Obligations of a process transition $\obl{\wrld ; \senv}{n}{\proc \trans{\pi} \proc' : \ty}$ (\rocq{Pobl})}
  \begin{proofrules}
    \infer*[lab=Q-Tau]{
      \proc'.1 \disjoint \dom(\senv) \implies
      \exists\, N \text{ fresh for } \senv, \ty,\
        \wrld \wextfresh \wrld' \text{ placing } N \st
      \proc' \in \lrPW{\wrld'}{n \cost{\pi}}{\senv \types \ty}
      \\
      \text{and, if the step is the $\tau$ of \ref{rule:P-Act-Link}},\
        \proc' \in \lrPW{\wrld}{n}{\senv \types \ty}
    }{
      \obl{\wrld ; \senv}{n}{\proc \trans{\tau} \proc' : \ty}
    } \label{rule:Q-Tau}
    \and
    \infer*[lab=Q-Recv]{
      \forall \wrld' \wextsym \wrld,\ j \le n \cost{\pi},\ \tenv_2
        \text{ with } \dom(\tenv_2) \disjoint \proc.1,\
        \val \in \lrVW{\wrld'}{j}{\tenv_2 \types \tysub{\tyB}{\pntA}{\pself}} \st
      \\
      (\proc'.1 \cup B, \proc'.2) \in
        \lrPW{\wrld'}{j}{(\senv_1 , \ch : \sAt{\sty}{(\poff{\pntA}{1})}) \mathbin{\triangleright} \tenv_2 \types \ty}
    }{
      \obl{\wrld ; \senv_1 , \ch : \sAt{(\sRecv{\tyB}{\sty})}{\pntA}}{n}{\proc \trans{\lRecv{\ch}{\val}} \proc' : \ty}
    } \label{rule:Q-Recv}
    \and
    \infer*[lab=Q-Send]{
      \dom(\tenv_{\asTy{e}}) = \fchans(\val) \cap \proc.1 \\
      \tenv_{\asTy{e}} \text{ and } \senv_{\asTy{e}} \text{ are dual at aliased points}
      \\\\
      \val \in \lrVW{\wrld}{n}{\senv_{\asTy{v}} , \tenv_{\asTy{e}} \types \tysub{\tyB}{\pntA}{\pself}}
      \quad
      \proc' \in \lrPW{\wrld}{n \cost{\pi}}{\senv_0 , \senv_{\asTy{e}} , \ch : \sAt{\sty}{(\poff{\pntA}{1})} \types \ty}
    }{
      \obl{\wrld ; \senv_0 , \senv_{\asTy{v}} , \ch : \sAt{(\sSend{\tyB}{\sty})}{\pntA}}{n}{\proc \trans{\lSend{\ch}{\val}} \proc' : \ty}
    } \label{rule:Q-Send}
    \and
    \infer*[lab=Q-Close]{
      \proc' \in \lrPW{\wrld}{n \cost{\pi}}{\senv_1 \types \ty}
    }{
      \obl{\wrld ; \senv_1 , \ch : \sAt{\sTerm}{\pntA}}{n}{\proc \trans{\lTerm{\ch}} \proc' : \ty}
    } \label{rule:Q-Close}
    \and
    \infer*[lab=Q-Link]{
      \ch \ne \chB \\
      \sty_{\chB} \isdual \sty_{\ch} \\
      \wrld \types \fEq{\pntA}{\pntB} \\
      \proc' \in \lrPW{\wrld}{n \cost{\pi}}{\senv_1 \types \ty}
    }{
      \obl{\wrld ; \senv_1 , \ch : \sAt{\sty_{\ch}}{\pntA} , \chB : \sAt{\sty_{\chB}}{\pntB}}{n}{\proc \trans{\lLink{\ch}{\chB}} \proc' : \ty}
    } \label{rule:Q-Link}
  \end{proofrules}
  \rulesectionend
  \caption{The process logical relation}
  \Description{The two rules of the process relation and the five
    obligation rules of a process transition.}
  \label{fig:logrel-rules-proc}
\end{figure}

\Cref{rule:Q-Tau} is the only clause that grows the world, by the fresh
extension that \ref{rule:P-Act-New} needs; a $\tau$ of
\ref{rule:P-Act-Link} creates no point and keeps the world, and it also
gives away a resource, so it carries its obligation at $n$ as well.
In \ref{rule:Q-Recv} the message may bring the other end of a channel
the receiver already holds, or both ends of a fresh one. Those channels
$B$ become restricted in the continuation and their entries leave both
environments, which the operation $\senv \mathbin{\triangleright} \tenv_2$
performs (\rocq{recv_back}, \rocq{recv_env}).
\Cref{rule:Q-Send} is where a private channel becomes part of the
external interface. The restricted channels the message carries are
$\dom(\tenv_{\asTy{e}}) = \fchans(\val) \cap \chset$; the ends that
leave with the message are typed in $\tenv_{\asTy{e}}$, the ends that
stay in $\senv_{\asTy{e}}$, and the two are dual at aliased points
(\rocq{send_split}). The continuation then serves
$\senv_{\asTy{e}}$ as part of its interface. No condition on the
topology of the continuation is imposed: the order on points does that
work, so this rule needs no split of the continuation into one process
per extruded channel.

The syntactic side condition $\fchans(\val)$ in \ref{rule:Q-Send} is
again a consequence of the design of the LTS (\cref{rm:sentc}).

\subsection{Fundamental Property}
\label{sec:fundamental}

The term relation applies to closed terms, so an open term is closed by
a grounding substitution.

\begin{definition}[Grounding substitutions \rocq{grounding}]
\label{def:grounding}
  \[
    \lrGW{\wrld}{n}{\tenv ; \vr_1 : \ty_1, \dots, \vr_k : \ty_k} :=
      \left\{
        [\vr_1/\val_1,\dots,\vr_k/\val_k] \;\middle|\;
          \begin{matrix}
          \tenv = \tenv_1 \cdot {} \dots \cdot \tenv_k
            \\
          \forall i.\,
            \val_i \in \lrVW{\wrld}{n}{\tenv_i \types \ty_i}
          \end{matrix}
      \right\}
  \]
\end{definition}

\begin{definition}[Semantic judgments \rocq{sem_typed}, \rocq{sem_ptyped}]
\label{def:sem-judgments}
  The judgment $\tenv ; \lenv \provesF{\frm}_{\wrld} \tm : \ty$ holds if
  \[
    \forall n \in \Nat.\,
      \forall \wrld' \wextsym \wrld.\,
        \forall \tenv'', \gamma \in \lrGW{\wrld'}{n}{\tenv' ; \lenv}
          \text{ with } \tenv'' = \tenv \cdot \tenv'.\,
          \tm\gamma \in \lrEW{\wrld'}{n}{\tenv'' \types \ty}{\frm}
  \]
  The judgment $\senv \proves_{\wrld} \proc : \ty$ holds if
  $
    \forall n \in \Nat.\,
      \proc \in \lrPW{\wrld}{n}{\senv \types \ty}
  $.
\end{definition}

The judgment is purely semantic. It says nothing about the syntax of
$\tm$ beyond what the term relation says about its behaviour. The two
syntactic conditions of \cref{rm:vabs} live in a separate predicate
$\mathit{wf}(\tenv; \lenv, \tm)$ (\rocq{term_wf}), which holds of every
syntactically typed term (\rocq{typed_term_wf_s}) and enters as a
premise of exactly the two compositionality lemmas that build a
suspended computation.

\begin{lemma}[\rocq{Prel_steps}]
\label{lm:lrproc-inv}
  If $\proc \in \lrPW{\wrld}{n+m}{\tEmpty \types \ty}$ and
  $\proc \redu[n] \proc'$, then
  $\proc' \in \lrPW{\wrld'}{m}{\tEmpty \types \ty}$ for some $\wrld'$.
\end{lemma}

\begin{lemma}[\rocq{Prel_safe}]
\label{lm:lrproc-safe}
  If $\proc \in \lrPW{\wrld}{n+1}{\tEmpty \types \ty}$ then $\proc$ is
  safe.
\end{lemma}

\begin{theorem}[Adequacy \rocq{adequacy}]
  \label{thm:adequacy}
  If $\tEmpty \proves_{\wrld} \proc : \ty$ then $\proc$ is deadlock
  free.
\end{theorem}

\begin{proof}
  Let $\proc \redu[m] \proc'$. Pick $n = m+1$ in
  $\tEmpty \proves_{\wrld} \proc : \ty$. By
  \cref{lm:lrproc-inv}, $\proc' \in \lrPW{\wrld'}{1}{\tEmpty \types \ty}$
  for some $\wrld'$, and by \cref{lm:lrproc-safe} $\proc'$ is safe.
\end{proof}

The relation validates every rule of the static system. Replacing
$\types$ by $\proves$ in \cref{fig:typeSys,fig:typeSysPoints,fig:procTypeSys} and reading
each inference rule as an implication gives a valid statement about the
semantic judgment, with the two syntactic premises noted above.
\Cref{tab:compositionality} names the declaration for each rule, and the
three lemmas below spell out the rules that carry a side condition of
their own.

\begin{table}[tp]
  \caption{Compositionality: one declaration per rule of the static
    system.}
  \label{tab:compositionality}
  \begin{tabular}{llllll}
    \toprule
    Rule & Declaration & Rule & Declaration & Rule & Declaration \\
    \midrule
    \ref{rule:E-Unit} & \rocq{sem_typed_unit}
      & \ref{rule:E-Const} & \rocq{sem_typed_konst}
      & \ref{rule:E-Var} & \rocq{sem_typed_var} \\
    \ref{rule:E-Abs} & \rocq{sem_typed_lam}
      & \ref{rule:E-App} & \rocq{sem_typed_app}
      & \ref{rule:E-Pair} & \rocq{sem_typed_pair} \\
    \ref{rule:E-LetPair} & \rocq{sem_typed_letp}
      & \ref{rule:E-Let*} & \rocq{sem_typed_letu}
      & \ref{rule:E-Inj} & \rocq{sem_typed_inj1} \\
    \ref{rule:E-Case} & \rocq{sem_typed_case}
      & \ref{rule:E-Weak} & \rocq{sem_typed_wk}
      & \ref{rule:E-Contract} & \rocq{sem_typed_contract} \\
    \ref{rule:E-Conv} & \rocq{sem_typed_teq}
      & \ref{rule:E-Exp} & \rocq{sem_typed_box}
      & \ref{rule:E-Pack} & \rocq{sem_typed_pack} \\
    \ref{rule:E-Open} & \rocq{sem_typed_open}
      & \ref{rule:P-Term} & \rocq{Prel_term}
      & \ref{rule:P-Connect} & \rocq{Prel_connect} \\
    \bottomrule
  \end{tabular}
\end{table}

A channel constant is typed by its own rule in the world-carrying
judgment, with declaration \rocq{sem_typed_chan}, and
$\tApp{\tForce}{\tm}$ is \ref{rule:E-App} at $\tForce$
(\rocq{sem_typed_force}). Composing two processes that do not interact
is \ref{rule:P-Connect} with no shared channel (\rocq{Prel_indep}).

\begin{lemma}[Compositionality: \ref{rule:E-Abs} \rocq{sem_typed_lam}]
  \label{lem:comp:abs}
  If $\tenv ; \lenv , \vr : \ty \provesF{\set{\pfv}}_{\wrld} \tm : \tyB$
  and $\wrld \types \fLe{\lbd}{\heads(\tenv) \cup \heads(\lenv)}$ and
  $\mathit{wf}(\tenv ; \lenv , \vr : \ty, \tm)$, then
  $\tenv ; \lenv \provesF{\frm}_{\wrld} \tAbs{\vr}{\tm}
    : \ty \lolli{\ubd}{\lbd} \tyB$.
\end{lemma}

\begin{lemma}[Compositionality: \ref{rule:E-Exp} \rocq{sem_typed_box}]
  \label{lem:comp:exp}
  If
  $\tEmpty ; \lenv , \vr : \tyBox{\pvivec}{\cstr}{\tyUnit \lolli{\ubd}{\pTop} \ty}
     \provesF{\set{\pfv}}_{\wrld} \tm : \ty$
  with $\lenv$ of $\Box$ types only and
  $\mathit{wf}(\tEmpty ; \lenv , \vr : \tyBox{\pvivec}{\cstr}{\tyUnit \lolli{\ubd}{\pTop} \ty}, \tm)$,
  then
  $\tEmpty ; \lenv \provesF{\frm}_{\wrld} \tBoxP{\vr}{\tm}
    : \tyBox{\pvivec}{\cstr}{\tyUnit \lolli{\ubd}{\pTop} \ty}$.
\end{lemma}

\begin{lemma}[Compositionality: \ref{rule:P-Connect} \rocq{Prel_connect}]
  \label{lem:comp:connect}
  If $\primitive \ty$, the two interfaces
  $\senv_{\asTy{C}}$ and $\dual{\senv_{\asTy{C}}}$ are dual at aliased
  points of $\wrld$,
  $(\chset_1, \pool_1) \in \lrPW{\wrld}{n}{\senv_1 , \senv_{\asTy{C}} \types \ty}$,
  $(\chset_2, \pool_2) \in \lrPW{\wrld}{n}{\senv_2 , \dual{\senv_{\asTy{C}}} \types \tyUnit}$,
  the restricted channels of each component avoid the interface of the
  other, and $\wrld \sqcup \wrld$ is strict, then
  \[
    (\chset_1 \cupdot \chset_2 \cupdot \dom(\senv_{\asTy{C}}),
     \pool_1 \uplus \pool_2)
      \in \lrPW{\wrld}{n}{\senv_1 , \senv_2 \types \ty}
  \]
\end{lemma}

\Cref{lem:comp:connect} is the linear-logic cut, and it is the one place
where strictness of the composed world is needed:
\cref{thm:amalgamation,thm:fresh-extension} supply it along the run,
because the only clause that grows a world grows it freshly
(\ref{rule:Q-Tau}) and the names it adds are unknown to the other
component.

\begin{theorem}[Fundamental property \rocq{fundamental}, \rocq{fundamental_proc}]
\label{th:fundamental}
  Typing implies semantic typing, at every strict world.
  \begin{itemize}
    \item
      $
        \fctx ; \lenv \typesF{\frm} \tm : \ty
        \implies
        \tenv ; \lenv \provesF{\frm}_{\wrld} \tm : \ty
      $
    \item
      $
        \senv \types \proc : \ty
        \implies
        \senv \proves_{\wrld} \proc : \ty
      $
  \end{itemize}
\end{theorem}

\begin{theorem}[Soundness \rocq{soundness}]
\label{th:soundness}
  If $\wrld$ is strict and $\tEmpty \types \proc : \ty$, then $\proc$ is
  deadlock free.
\end{theorem}

\Cref{thm:soundness-source} is \cref{th:soundness} for a source program:
the judgment of \cref{fig:typeSys} and \cref{fig:typeSysPoints} reads no world, the acyclic
$\cstr_0$ supplies a strict world that realizes it, and
\rocq{styped_typed} embeds the derivation (\rocq{soundness_source}).

%%% Local Variables:
%%% mode: latex
%%% TeX-master: "main"
%%% End:
